\section{A sparse prime-sum criterion for simple zeros}
\label{sec:sparse-simplicity}

The logarithmic-derivative method behind the Lean reduction also detects
multiplicity. For the Riemann zeta function it gives the following precise
target. Let $\gamma$ denote Euler's constant and $\gamma_1$ the first Stieltjes
constant, with $\zeta(1+u)=u^{-1}+\gamma-\gamma_1u+O(u^2)$. Set
\[
 x_n=2^n3^{n^2},\qquad
 R_{\mathbb Q}(x)=
 \sum_{p^j\le x,\ j\ge1}\log p\left(\frac{1}{p^j}-\frac1x\right)
 -\log x+1+\gamma .
\]
\noindent\textbf{Simplicity criterion.}
Under RH, all nontrivial zeros of $\zeta$ are simple if and only if
\[
 \lim_{N\to\infty}\frac1N\sum_{n=1}^N x_nR_{\mathbb Q}(x_n)^2
 =6+2\gamma-2\log(4\pi)+4\gamma_1+2\gamma^2-\frac{\pi^2}{4}.
\]
The constant is approximately $0.00007420127287493$. The limit always
exists under RH. Its excess above the displayed value is positive if any
zero is multiple. The remaining arithmetic question is an upper bound on
the sampled mean square by this explicit constant. Only
$O(\sqrt{\log X})$ cutoffs occur through $X$, although each cutoff
uses prime powers up to its full value. Continuous mean squares of prime
errors also detect repeated zeros. Here the reciprocal-prime-power
weight makes the Fourier series absolutely convergent, and the polynomial
choice of logarithms permits sampling on a single sparse sequence.
\footnote{For the continuous mean square of the unsmoothed Chebyshev error
and its dependence on multiplicities, see P.~Caceres,
\emph{The constant term in Cramér's mean-square theorem for the primes},
\href{https://hexagonmath.org/2611.01010v1}{preprint, 2026}.}

\paragraph{The zero expansion on the critical line.}
Use the same $R_K$ and explicit formula as in Section~1. In Ihara's notation
the uncentered prime sum equals $(1-x^{-1})\Phi_K(x)$.
The normalization and archive sources are recorded in \texttt{sources.md}.

Assume all nontrivial zeros of $\zeta_K$ lie on the critical line. For each
real $t$ let $m_K(t)$ be the multiplicity of $1/2+it$, with value zero at
nonzeros, and let $\Gamma_K$ be the set of distinct ordinates. Then
\[
 \sqrt{x}R_K(x)=\sum_{t\in\Gamma_K} a_t x^{it}+o(1),
 \qquad a_t=\frac{m_K(t)}{1/4+t^2},\qquad \sum_t a_t<\infty .
\]
Absolute summability gives uniform convergence in $\log x$. The Lean
repository supplies this line hypothesis when $K=\mathbb Q$. For abelian
extensions of $\mathbb Q$ it follows from the factorization into Dirichlet
$L$-functions. For a general number field it is an additional hypothesis.

\paragraph{Recovery of every multiplicity.}
For every real $u$,
\[
 \lim_{N\to\infty}\frac1N\sum_{n=1}^N x_n^{iu}
   =\begin{cases}1&u=0,\\0&u\ne0.\end{cases}
\]
Indeed, for $u\ne0$, at least one of $u\log2/(2\pi)$ and $u\log3/(2\pi)$ is
irrational, since $\log2/\log3$ is irrational by unique factorization. Apply
Weyl's theorem to the polynomial with these linear and quadratic
coefficients. The same argument works with $n^2$ replaced by $n^d$ for
each fixed integer $d\ge2$.
\footnote{H.~Weyl, \emph{\"Uber die Gleichverteilung von Zahlen mod.\ Eins},
Math. Ann. 77 (1916), 313--352, \href{https://doi.org/10.1007/BF01475864}{doi:10.1007/BF01475864}.}
Multiplication by $x_n^{-iu}$ and termwise averaging now give
\[
 \lim_{N\to\infty}\frac1N\sum_{n=1}^N
 \sqrt{x_n}R_K(x_n)x_n^{-iu}=\frac{m_K(u)}{1/4+u^2}.
\]
This recovery formula requires no assumptions about rational relations
among zero ordinates. Absolute convergence controls the tail uniformly in
$N$, and the average of the $o(1)$ term tends to zero.

\paragraph{The defect from simplicity.}
Since the prime sum is real, its square is its squared absolute value.
Applying the same averaging identity to pairs of ordinates gives
\[
 \lim_{N\to\infty}\frac1N\sum_{n=1}^N x_nR_K(x_n)^2
     =\sum_{t\in\Gamma_K}\frac{m_K(t)^2}{(1/4+t^2)^2}.
\]
Indeed the double sum is absolutely convergent, and only terms whose
ordinates coincide survive. Write $H_K=(\log\xi_K)'$, where
$\xi_K(s)=s(s-1)|D_K|^{s/2}\Gamma_{\mathbb R}(s)^{r_1}
\Gamma_{\mathbb C}(s)^{r_2}\zeta_K(s)$ is the completed entire zeta
function, $D_K$ is the discriminant and $(r_1,r_2)$ the signature.
Here $\Gamma_{\mathbb R}(s)=\pi^{-s/2}\Gamma(s/2)$ and
$\Gamma_{\mathbb C}(s)=2(2\pi)^{-s}\Gamma(s)$.
The genus-one Hadamard product gives
\[
 \sum_\rho\frac1{\rho(1-\rho)}
       =H_K(1)-H_K(0)=2H_K(1),\qquad
 \sum_\rho\frac1{(1-\rho)^2}=-H'_K(1).
\]
The first identity uses $\xi_K(s)=\xi_K(1-s)$.
The identity
$[z(1-z)]^{-2}=z^{-2}+(1-z)^{-2}+2/[z(1-z)]$ and the same symmetry give
\[
 \sum_{t\in\Gamma_K}\frac{m_K(t)}{(1/4+t^2)^2}
      =4H_K(1)-2H'_K(1).
\]
Consequently the difference between the sampled mean square and
$4H_K(1)-2H'_K(1)$ is exactly
$\sum_t m_K(t)(m_K(t)-1)/(1/4+t^2)^2$. Every term is nonnegative.
For $K=\mathbb Q$ the expansion at $s=1$ gives
\[
 H_{\mathbb Q}(1)=1+\tfrac12\gamma-\tfrac12\log(4\pi),\qquad
 H'_{\mathbb Q}(1)=-1-2\gamma_1-\gamma^2+\pi^2/8,
\]
which proves the stated constant. Controlling the off-diagonal terms when
one expands the \emph{prime} side of the mean square at these cutoffs
remains open.

\paragraph{A numerical enclosure for the limiting statistic.}
Let $D$ be the excess in the Riemann case and $H=3\cdot10^{12}$.
Then
\[
 0\le D <1.67\cdot10^{-34}.
\]
Platt and Trudgian's disjoint sign changes exhaust the zero count with multiplicity through $H$, so each of those zeros is simple. For $t\ge H$ Corollary~1.2 of Hasanalizade, Shen and Wong
bounds the difference of the zero count $N(t)$ from
$(t/2\pi)\log(t/(2\pi e))$ by
$0.1038\log t+0.2573\log\log t+9.3675<\log t$.
Hence $N(t)\le t\log t$ there, and a jump at $t>H$ has size at most
$2\log t$. By conjugation symmetry and Stieltjes integration by parts,
\[
 D\le 4\int_{(H,\infty)}\frac{\log t}{t^4}\,dN(t)
 \le16\int_H^\infty\frac{\log^2t}{t^4}\,dt
 =\frac{16}{H^3}\left(\frac{\log^2H}3+
     \frac{2\log H}9+\frac2{27}\right).
\]
The boundary term discarded at $H$ is nonpositive. A multiple zero at
arbitrarily large height could contribute less than this bound.
\footnote{D.~Platt and T.~Trudgian,
\emph{The Riemann hypothesis is true up to $3\cdot10^{12}$},
\href{https://doi.org/10.1112/blms.12460}{Bull. London Math. Soc. 53 (2021), 792--797}.
E.~Hasanalizade, Q.~Shen and P.-J.~Wong,
\href{https://arxiv.org/abs/2107.06506}{\emph{Counting zeros of the Riemann zeta function}},
Corollary~1.2. The earlier ProveIt report on Suzuki's density question uses
the same two estimates with a different spectral weight.}
